\documentclass[11pt,letterpaper]{article}

\usepackage[top=1in, bottom=1in, left=1in, right=1in, headheight=14pt]{geometry}
\usepackage{fontspec}
\setmainfont{DejaVu Serif}
\setsansfont{DejaVu Sans}
\setmonofont{DejaVu Sans Mono}

\usepackage{xcolor}
\usepackage{titlesec}
\usepackage{fancyhdr}
\usepackage{enumitem}
\usepackage{hyperref}
\usepackage{booktabs}
\usepackage{tabularx}
\usepackage{longtable}
\usepackage{colortbl}
\usepackage{multirow}
\usepackage{array}
\usepackage{parskip}
\usepackage{tcolorbox}
\usepackage{graphicx}
\usepackage{lastpage}

%% Color Scheme -- Volcanic / Governance
\definecolor{primary}{HTML}{5D1A00}       % Volcanic red-brown — Section headers
\definecolor{secondary}{HTML}{0D3B5E}     % Deep ocean blue — Subsection headers
\definecolor{tertiary}{HTML}{7B5E00}      % Ember gold — Subsubsection headers
\definecolor{accent}{HTML}{8C2E00}        % Lava orange-red — Highlights
\definecolor{lightprimary}{HTML}{FBE9E7}  % Light volcanic rose — Quote boxes
\definecolor{lightsecondary}{HTML}{E3EFF7}
\definecolor{lighttertiary}{HTML}{FFF8E1}
\definecolor{rowgray}{HTML}{F5F5F5}
\definecolor{bordergray}{HTML}{BDBDBD}
\definecolor{subtlegray}{HTML}{757575}
\definecolor{darktext}{HTML}{212121}

%% Section formatting
\titleformat{\section}
  {\Large\bfseries\sffamily\color{primary}}
  {\thesection}{1em}{}
  [\vspace{-0.5em}\textcolor{bordergray}{\rule{\textwidth}{0.4pt}}]

\titleformat{\subsection}
  {\large\bfseries\sffamily\color{secondary}}
  {\thesubsection}{1em}{}

\titleformat{\subsubsection}
  {\normalsize\bfseries\sffamily\color{tertiary}}
  {\thesubsubsection}{1em}{}

\titlespacing*{\section}{0pt}{1.5em}{0.8em}
\titlespacing*{\subsection}{0pt}{1.2em}{0.5em}
\titlespacing*{\subsubsection}{0pt}{0.8em}{0.3em}

%% Custom environments
\newcommand{\stageheader}[2]{%
  \noindent\colorbox{#2}{\makebox[\textwidth][l]{%
    \textcolor{white}{\bfseries\sffamily\hspace{6pt}#1\hspace{6pt}}}}%
  \vspace{0.5em}\par%
}

\newtcolorbox{asciibox}{
  colback=rowgray, colframe=bordergray,
  arc=1pt, boxrule=0.5pt,
  left=6pt, right=6pt, top=4pt, bottom=4pt,
  fontupper=\small\ttfamily,
  breakable
}

\newtcolorbox{quotebox}{
  colback=lightprimary, colframe=primary,
  arc=2pt, boxrule=0.5pt,
  left=12pt, right=12pt, top=8pt, bottom=8pt,
  breakable
}

\newcommand{\metafield}[2]{\textbf{\sffamily #1} #2\par}
\newcolumntype{L}[1]{>{\raggedright\arraybackslash}p{#1}}
\newcolumntype{C}[1]{>{\centering\arraybackslash}p{#1}}
\newcommand{\tableheader}[1]{\textbf{\sffamily\textcolor{white}{#1}}}

%% Headers and Footers
\pagestyle{fancy}
\fancyhf{}
\fancyhead[L]{\small\sffamily\textcolor{subtlegray}{Ring of Fire: Tribal Sovereignty \& Governance}}
\fancyhead[R]{\small\sffamily\textcolor{subtlegray}{GSD Mission Package}}
\fancyfoot[C]{\small\sffamily\textcolor{subtlegray}{Page \thepage\ of \pageref{LastPage}}}
\renewcommand{\headrulewidth}{0.4pt}
\renewcommand{\footrulewidth}{0pt}

\hypersetup{
  colorlinks=true,
  linkcolor=secondary,
  urlcolor=secondary,
  pdfauthor={GSD Ecosystem / Tibsfox Inc.},
  pdftitle={Ring of Fire: Tribal Sovereignty and Governance -- Mission Package},
  pdfsubject={Vision to Mission Pipeline},
}

\begin{document}

%% ============================================================
%% TITLE PAGE
%% ============================================================
\thispagestyle{empty}
\begin{center}
\vspace*{3cm}

{\small\sffamily\textcolor{primary}{\textbf{GSD RESEARCH MISSION PACKAGE}}}

\vspace{0.3cm}
\textcolor{bordergray}{\rule{8cm}{0.4pt}}

\vspace{1.5cm}
{\fontsize{28}{34}\selectfont\bfseries\sffamily\textcolor{primary}{Ring of Fire\\[0.3em]Tribal Sovereignty \& Governance}}

\vspace{0.8cm}
{\Large\sffamily\textcolor{secondary}{Pacific Rim — A 40,000 km Arc of Living Nations}}

\vspace{0.5cm}
{\large\sffamily\itshape\textcolor{tertiary}{A Deep Research Study}}

\vspace{3cm}
{\sffamily\textcolor{accent}{Vision $\rightarrow$ Research $\rightarrow$ Mission}}\\[0.3em]
{\small\sffamily\textcolor{accent}{Full Pipeline Execution}}

\vspace{3cm}
{\small\sffamily\textcolor{subtlegray}{Date: April 2026  |  Status: Ready for GSD Orchestrator}}\\[0.3em]
{\small\sffamily\textcolor{subtlegray}{Activation Profile: Squadron (12 roles)  |  Estimated: 18--24 context windows across 4--6 sessions}}
\end{center}
\newpage

%% ============================================================
%% TABLE OF CONTENTS
%% ============================================================
\tableofcontents
\newpage

%% ============================================================
%% STAGE 1 — VISION DOCUMENT
%% ============================================================
\stageheader{STAGE 1 --- VISION DOCUMENT}{primary}

\section{Ring of Fire: Tribal Sovereignty \& Governance --- Vision Guide}

\metafield{Date:}{April 2026}
\metafield{Status:}{Ready for GSD Orchestrator — Full Pipeline Execution}
\metafield{Depends On:}{Fox Infrastructure Group Plan Review, Tibsfox Inc. Articles of Incorporation}
\metafield{Context:}{Pacific Spine corridor (BC to Chile); Fox Infrastructure Group; Ring of Fire / Fox Infrastructure Group vision}
\metafield{Activation Profile:}{Squadron (12 roles, 5 research modules)}

\subsection{Vision}

There is a horseshoe of fire drawn around the Pacific Ocean. It runs 40,000 kilometers --- from the tip of South America, up through Chile, Peru, Mexico, California, Oregon, Washington, British Columbia, Alaska, across the Aleutian chain, down through Japan, the Philippines, Indonesia, Papua New Guinea, and New Zealand. It is the zone where tectonic plates collide and subduct, producing ninety percent of the world's earthquakes and the majority of its active volcanoes.

But the Ring of Fire is not simply a geological phenomenon. It is also the longest continuous corridor of Indigenous nations on the planet. Every kilometer of that arc has been lived on, stewarded, governed, and named by peoples whose presence predates any nation-state boundary drawn across the Pacific. The Quechua and Aymara of the Andes. The coastal nations of Mexico and California --- the Ohlone, the Chumash, the Tongva. The Coast Salish peoples of the Salish Sea --- the Lummi, Tulalip, Stillaguamish, Sauk-Suiattle, Suquamish, Snohomish. The Ainu of Hokkaido and the Kuril Islands. The Māori of Aotearoa New Zealand. The Igorot, Lumad, and Moro peoples of the Philippines.

What is now called the Ring of Fire was, before colonization, a Ring of Governance.

The Fox Infrastructure Group envisions a Pacific Spine corridor running from British Columbia to Chile --- a cooperative infrastructure network rooted in communities, not extraction. Any such network that ignores the sovereign nations along its path is not a network; it is a pipeline drawn through other people's living rooms. This research mission exists to build the foundational knowledge base for a partnership model that begins with sovereignty, not extraction. The question is not ``how do we get Indigenous consent?'' --- the UNDRIP framework answers that. The question is: \textit{how do we build infrastructure that Indigenous nations would want to own?}

The GSD ecosystem's prime directive --- giving people their lives back --- applies here at civilizational scale. Nations that have had their governance structures dismantled, their land bases fractured, and their economies deliberately impoverished are the communities most in need of the kind of friction-reducing, dignifying infrastructure this project aims to build. The Ring of Fire is the most important geopolitical and ecological corridor on Earth. Understanding its sovereign nations is not a compliance exercise; it is the foundational act of building anything here worth building.

\subsection{Problem Statement}

\begin{enumerate}
  \item \textbf{The consent gap:} Critical mineral extraction, infrastructure development, and energy projects across the Ring of Fire continue to proceed without genuine Free, Prior, and Informed Consent (FPIC) from affected Indigenous nations, creating legal uncertainty, project delays, conflict, and ongoing rights violations. The Canadian Ring of Fire in northern Ontario is the clearest current example, where 15 First Nations are engaged in a federally co-led Regional Assessment while development pressure mounts.

  \item \textbf{The fragmentation problem:} Pacific Rim Indigenous nations govern within dozens of different national legal frameworks --- from the Ainu Policy Promotion Act of 2019 in Japan, to the Treaty of Waitangi in New Zealand, to U.S. federal Indian law and Canadian treaty frameworks. There is no shared mapping of these governance structures that infrastructure planners, cooperative networks, or sovereign partnership frameworks can reliably use.

  \item \textbf{The extraction default:} The dominant model for engaging Indigenous nations along resource corridors is still extractive: consult minimally, share benefits nominally, retain control entirely. The emerging co-investment and equity-partnership models (Marten Falls First Nation equity in the Northern Road Link; Keewaytinook Okimakanak's KNET ISP; Wasaya Airways coalition ownership) represent viable alternatives but are not yet systematized.

  \item \textbf{Data and digital sovereignty lag:} Indigenous nations are asserting data sovereignty --- the right to govern the collection, ownership, and application of data about their communities --- but the infrastructure enabling that sovereignty (Indigenous-owned ISPs, data centers, digital governance frameworks) is nascent and unevenly distributed across the Pacific Rim.

  \item \textbf{Governance continuity under climate stress:} The Ring of Fire is both a seismic hazard zone and one of the world's most climate-vulnerable corridors. More than 50,000 Pacific Island peoples are displaced annually by climate events. Traditional governance structures encode ecological knowledge essential for resilience, but are not integrated into disaster preparedness or infrastructure planning frameworks.
\end{enumerate}

\subsection{Core Concept}

\textbf{The Sovereignty Stack:} Governance, Territory, Resources, Data, and Infrastructure are five vertically nested layers of sovereignty. No infrastructure project along the Ring of Fire is viable unless it engages all five layers --- beginning at the top.

\begin{asciibox}
SOVEREIGNTY STACK — Ring of Fire Governance Framework

Layer 5: INFRASTRUCTURE
    Physical networks, energy, fiber, transport, compute
    "Who owns the pipe?"

Layer 4: DATA SOVEREIGNTY
    Collection, ownership, application of community data
    "Who controls the signal?"

Layer 3: RESOURCE RIGHTS
    Minerals, fisheries, water, geothermal, forest
    "Who owns what is extracted?"

Layer 2: TERRITORIAL AUTHORITY
    Land tenure, treaty boundaries, traditional use areas
    "Whose land is this?"

Layer 1: GOVERNANCE FOUNDATION
    Inherent sovereignty, constitutions, self-determination
    "Who makes decisions here?"

Principle: You cannot build Layer 5 without engaging Layer 1.
FPIC is not a gate — it is the foundation material.
\end{asciibox}

\subsubsection{Pacific Rim Arc Map}

\begin{asciibox}
RING OF FIRE — INDIGENOUS SOVEREIGNTY ARC (Selected Nations)

    Alaska / Aleutians
    [Unangan / Aleut Nations]
           |
    Pacific Northwest / Salish Sea
    [Lummi, Tulalip, Stillaguamish, Sauk-Suiattle,
     Suquamish, Makah, Quinault, Coast Salish (BC)]
           |
    California / Oregon
    [Karuk, Yurok, Hupa, Ohlone, Chumash]
           |
    Mexico / Central America
    [Nahua, Zapotec, Maya coastal nations]
           |
    South America (Andean Spine)
    [Quechua, Aymara, Mapuche, Rapa Nui]

    =========== PACIFIC ============

    Japan / Hokkaido / Kurils
    [Ainu (recognized 2008), Ryukyuans (Okinawa, unrecognized)]
           |
    Philippines
    [Igorot, Lumad, Moro, Mangyan]
           |
    Indonesia / Papua New Guinea
    [Papuan nations, Maluku, West Papua]
           |
    Aotearoa New Zealand
    [Maori (Te Tiriti o Waitangi, 1840)]

All nations: Inherent sovereignty predating state formation.
All nations: Ongoing engagement with UNDRIP framework.
\end{asciibox}

\subsubsection{Module Architecture}

\begin{asciibox}
RESEARCH MODULE ARCHITECTURE

Module A: Territorial Rights & Legal Frameworks
    UNDRIP / FPIC / ILO 169 / National implementations
    [Sonnet — Survey]

Module B: Nation-by-Nation Governance Profiles
    PNW Coast Salish, Ainu, Maori, Andean, Philippines
    [Opus — Judgment-heavy, nation-specific]

Module C: Resource Sovereignty & Critical Minerals
    Ring of Fire (Ontario), Pacific geothermal, fisheries
    [Sonnet — Data compilation]

Module D: Cooperative Infrastructure Models
    Equity co-investment, ISPs, FoxInfrastructure integration
    [Opus — Synthesis, pattern extraction]

Module E: Indigenous Data & Digital Sovereignty
    CARE/OCAP, data governance frameworks, digital infra
    [Sonnet — Framework survey]

Wave 1: A + C (parallel)
Wave 1: B + E (parallel)
Wave 2: D (synthesis — depends on A, B, C, E)
\end{asciibox}

\subsection{Research Modules}

\subsubsection{Module A: Territorial Rights and Legal Frameworks}

Survey the international and national legal architecture governing Indigenous sovereignty along the Ring of Fire. Core instruments include the United Nations Declaration on the Rights of Indigenous Peoples (UNDRIP, 2007), ILO Convention 169, and national implementations ranging from Canada's UNDRIP Act (2021) to New Zealand's Treaty of Waitangi, Japan's Ainu Policy Promotion Act (2019), and U.S. federal Indian law. Particular attention to FPIC doctrine: its evolution from consultation requirement to consent requirement, the gap between endorsement and implementation, and the emerging equity co-investment model as the strongest FPIC signal.

\subsubsection{Module B: Nation-by-Nation Governance Profiles}

Produce authoritative, nation-specific governance profiles for the primary peoples along the Fox Infrastructure Group's Pacific Spine corridor. Profiles must use nation-specific names --- never generic identifiers. Priority nations: Lummi (Lhaq'temish), Tulalip, Stillaguamish, Sauk-Suiattle, Suquamish (PNW/Salish Sea); Neskantaga, Nishnawbe Aski Nation constituent communities (Ontario Ring of Fire); Ainu of Hokkaido (Japan); Māori iwi of the North Island volcanic arc (Aotearoa); Mapuche of southern Chile; Quechua/Aymara of the Andean altiplano. Each profile covers: governance structure, treaty/legal status, economic development posture, known infrastructure positions, key contacts and organizations.

\subsubsection{Module C: Resource Sovereignty and Critical Minerals}

Map the resource sovereignty conflicts and partnership models across the Ring of Fire. Focus areas: the Canadian Ring of Fire (James Bay Lowlands --- chromite, cobalt, nickel, copper, platinum; 15 First Nations in Regional Assessment; Marten Falls First Nation equity model); Pacific geothermal potential (the Ring holds over 40\% of world geothermal capacity --- Philippines, Indonesia, New Zealand, USA are leaders; Indigenous land tenure conflicts and partnership models); Pacific fisheries (Coast Salish reef net rights, Boldt Decision 1974, salmon sovereignty as infrastructure precondition). Document the shift from ``sacrifice zones'' to co-investor models.

\subsubsection{Module D: Cooperative Infrastructure Models}

Extract and systematize the emerging cooperative infrastructure models that have succeeded along the Ring of Fire. Cases include: Keewaytinook Okimakanak's KNET (First Nations-led ISP, now expanding into drone mapping and AI); Wasaya Airways (coalition First Nations ownership, Ring of Fire logistics); Lummi Nation's successful opposition to the Gateway Pacific Terminal at Cherry Point (treaty rights as infrastructure veto power); Māori co-governance of waterways under the Whanganui River personhood framework; Ainu land/sea rights advocacy using UN human rights instruments. Connect findings to Fox Infrastructure Group's co-op subsidiary model (FoxCompute, FoxFiber, FoxData, FoxEdu, FoxEnergy) and tribal sovereign partnership architecture.

\subsubsection{Module E: Indigenous Data and Digital Sovereignty}

Survey the Indigenous data sovereignty movement across the Pacific Rim, including: the CARE Principles (Collective Benefit, Authority to Control, Responsibility, Ethics); OCAP\textsuperscript{\textregistered} (Ownership, Control, Access, Possession --- First Nations Information Governance Centre, Canada); Māori data sovereignty (tino-rangatiratanga, Te Tiriti, kaupapa Māori data governance); the U.S. Indigenous Data Sovereignty Network; the Center for Tribal Digital Sovereignty (NCAI / Arizona State). Document emerging Indigenous-owned digital infrastructure: tribal ISPs, data centers, AI governance protocols. Connect to GSD ecosystem's DACP (Deterministic Agent Communication Protocol) and the planning-bridge MCP server as a model for Indigenous-controlled data flows.

\subsection{Chipset Configuration}

\begin{asciibox}
chipset_config:
  mission: rim_fire_sovereignty_governance
  profile: squadron
  model_allocation:
    opus_30pct:
      - Module B: nation-by-nation governance profiles
        (judgment-heavy, cultural sensitivity, nation-specific)
      - Module D: cooperative infrastructure synthesis
        (pattern extraction across cases, FIG integration)
      - Wave 2: cross-module synthesis
    sonnet_60pct:
      - Module A: legal frameworks survey
      - Module C: resource sovereignty data compilation
      - Module E: data sovereignty framework survey
      - Wave 3: publication assembly, verification
    haiku_10pct:
      - Wave 0: shared schemas, source index templates
      - Nation profile templates (Module B scaffold)
      - Bibliography management

  sensitivity_flags:
    - OCAP_CARE: nation-specific attribution REQUIRED
    - sacred_knowledge: ABSOLUTE gate — no reproduction
    - FPIC_language: use exact UNDRIP Article citations
    - Ainu: use "Ainu" not "aboriginal"; Ryukyuans not recognized by Japan — note carefully
    - Maori: use macrons correctly (Māori, Aotearoa, etc.)
    - Coast_Salish: use nation-specific names (Lhaq'temish, not generic)
\end{asciibox}

\subsection{Scope Boundaries}

\subsubsection{In Scope (v1.0)}

\begin{itemize}
  \item International legal frameworks: UNDRIP, ILO 169, FPIC doctrine, national implementations
  \item Governance profiles for nations along the BC-to-Chile Pacific Spine and the Japan-to-New Zealand arc
  \item The Canadian Ring of Fire (James Bay Lowlands) as the primary current sovereignty conflict case
  \item Cooperative infrastructure and equity co-investment models along the arc
  \item Indigenous data sovereignty frameworks (CARE, OCAP\textsuperscript{\textregistered}, Māori data sovereignty)
  \item Connection to Fox Infrastructure Group co-op subsidiary architecture
  \item PNW Coast Salish nations (Snohomish County / Paine Field prototype hub context)
\end{itemize}

\subsubsection{Out of Scope (v1.0)}

\begin{itemize}
  \item Detailed nation profiles for nations outside the Pacific Spine corridor and Japan-NZ arc
  \item Proprietary Impact Benefit Agreement terms or confidential tribal governance documents
  \item Sacred knowledge, ceremonial content, or restricted cultural information of any nation
  \item Detailed financial modeling for specific FIG co-investment proposals (deferred to FIG mission)
  \item Nations of the Western Pacific interior (non-coastal Ring of Fire nations)
\end{itemize}

\subsection{Success Criteria}

\begin{enumerate}
  \item A complete legal framework survey covering UNDRIP, ILO 169, and national implementations in at least six Ring of Fire countries (Canada, USA, Japan, New Zealand, Philippines, Chile), all claims attributed to government or peer-reviewed sources.
  \item Nation-specific governance profiles for at least 12 distinct nations along the Pacific Spine and Japan-NZ arc, each using nation-specific names and correctly describing governance structure, legal status, and economic development posture.
  \item A documented case study of the Canadian Ring of Fire Regional Assessment (2025--ongoing), covering the 15 First Nations partnership, terms of reference, FPIC requirements, and equity co-investment models.
  \item A structured comparison of at least five operative cooperative infrastructure models (KNET, Wasaya, Lummi-Cherry Point, Whanganui River, Marten Falls equity model), each with documented outcomes.
  \item A synthesis of the Indigenous data sovereignty landscape across the Pacific Rim, covering CARE Principles, OCAP\textsuperscript{\textregistered}, Māori data sovereignty, and at least two emerging Indigenous-owned digital infrastructure examples.
  \item A FIG integration document mapping the Sovereignty Stack (five layers) to Fox Infrastructure Group's co-op subsidiary architecture (FoxFiber, FoxCompute, FoxData, FoxEnergy, FoxEdu), with identified partnership entry points per layer.
  \item All nation-specific content reviewed against the OCAP\textsuperscript{\textregistered}/CARE sensitivity checklist: zero use of generic identifiers, zero reproduction of sacred or restricted content.
  \item A dependency graph connecting Modules A--E, identifying which findings from each module are load-bearing for Module D synthesis.
  \item A verified bibliography of at least 25 sources, all government agencies, peer-reviewed research, or professional organizations --- zero entertainment media or unsourced claims.
  \item A final compiled PDF passing three-pass XeLaTeX compilation with zero errors, complete table of contents, and page-accurate footers.
  \item All FPIC language in deliverables uses exact UNDRIP Article citations (Articles 10, 19, 28, 32) rather than paraphrase.
  \item The mission package is ready for handoff to the GSD orchestrator with zero additional context required from Tibsfox.
\end{enumerate}

\subsection{Relationship to Other GSD Documents}

\begin{center}
\rowcolors{2}{rowgray}{white}
\begin{tabularx}{\textwidth}{L{4cm}L{3cm}X}
\toprule
\rowcolor{primary}
\tableheader{Document} & \tableheader{Relationship} & \tableheader{Notes} \\
\midrule
Fox Infrastructure Group Plan Review & Depends On & Tribal sovereign partnerships are FIG's primary engagement model \\
Tibsfox Inc. Articles of Incorporation & Depends On & Public benefit purpose clause anchors sovereignty-first approach \\
GSD Chipset Architecture Vision & Informs & Model assignment (Opus for judgment/cultural sensitivity) \\
Seattle 360 Geo (seattle\_360\_geo.csv) & Adjacent & Coast Salish nations are central to PNW cultural ecosystem \\
College of Knowledge (.college/) & Future Dependency & Nation governance profiles feed Indigenous studies curriculum \\
gsd-skill-creator dev branch & Target & Mission package handed to GSD orchestrator for execution \\
\bottomrule
\end{tabularx}
\end{center}

\subsection{The Through-Line}

The Ring of Fire burns in both directions. Geologically, it is energy released from the collision of massive, slow-moving forces --- energy that has built islands, raised mountains, and occasionally destroyed cities. Politically, it is the same. The collision between inherited colonial state structures and the inherent sovereignty of the nations that preceded them is releasing energy that has been stored for generations. The question is whether that energy produces destruction or transformation.

The Amiga Principle teaches that specialized execution paths produce outsized outcomes over brute-force resource application. The lesson for sovereign partnership is identical: a model that begins with genuine governance engagement --- that treats each nation as the specialized, irreplaceable authority on their own territory --- will outperform any model that attempts to brute-force consent through consultation theater. Not just ethically. Economically. Projects with genuine Indigenous co-investment go forward. Projects without it go to court.

The Fox Infrastructure Group is not building infrastructure \textit{for} the communities of the Pacific Rim. It is building infrastructure \textit{with} the peoples who have governed the Pacific Rim for millennia and who will govern it for millennia more. The Sovereignty Stack is the architectural blueprint. The Ring of Fire nations are the chipset. Everything else is peripheral.

\newpage

%% ============================================================
%% STAGE 2 — RESEARCH REFERENCE
%% ============================================================
\stageheader{STAGE 2 --- RESEARCH REFERENCE}{secondary}

\section{Ring of Fire: Tribal Sovereignty \& Governance --- Research Reference}

\subsection{How to Use This Document}

This reference compiles verified findings, organized by research module, for handoff to GSD subagents. Each subsection corresponds to one module. All claims are sourced. Nation-specific naming is used throughout.

\textbf{Key Source Organizations:}

\begin{itemize}
  \item United Nations Permanent Forum on Indigenous Issues (UNPFII)
  \item First Nations Information Governance Centre (FNIGC) --- OCAP\textsuperscript{\textregistered}
  \item Nishnawbe Aski Nation (NAN), northern Ontario
  \item Impact Assessment Agency of Canada (IAAC)
  \item Puget Sound Regional Council (PSRC)
  \item Snohomish County Tribal Coordination Element (2023)
  \item International Work Group for Indigenous Affairs (IWGIA)
  \item International Institute for Democracy and Electoral Assistance (IDEA)
  \item Karuk Tribe (Good Fire reports I and II, 2021, 2024)
  \item U.S. Indigenous Data Sovereignty Network (USIDSN)
  \item Native Nations Institute, University of Arizona
\end{itemize}

\subsection{Module A Reference: Territorial Rights and Legal Frameworks}

\subsubsection{UNDRIP and the FPIC Standard}

The United Nations Declaration on the Rights of Indigenous Peoples (2007) received 144 votes in favor, with the United States, Canada, Australia, and New Zealand --- all mineral-rich nations with significant Indigenous populations --- originally voting against. All four have since endorsed UNDRIP to varying degrees. Canada adopted UNDRIP into domestic law via the United Nations Declaration on the Rights of Indigenous Peoples Act in 2021.

Key UNDRIP articles governing infrastructure and development:

\begin{center}
\rowcolors{2}{rowgray}{white}
\begin{tabularx}{\textwidth}{C{2cm}X}
\toprule
\rowcolor{secondary}
\tableheader{Article} & \tableheader{Provision} \\
\midrule
Article 10 & No relocation without FPIC and just compensation \\
Article 11 & Right to maintain and protect cultural heritage and traditional knowledge \\
Article 19 & States shall consult and cooperate to obtain FPIC before implementing measures affecting Indigenous peoples \\
Article 28 & Right to redress for lands taken without FPIC \\
Article 32 & States shall consult to obtain FPIC before approving projects affecting lands or resources \\
\bottomrule
\end{tabularx}
\end{center}

\subsubsection{FPIC in Practice: From Consultation to Consent}

The legal standard has evolved significantly. Consultation --- the weaker standard --- requires informing and hearing Indigenous communities. Consent --- the stronger UNDRIP standard --- requires agreement before proceeding. The distinction matters enormously in project outcomes: projects proceeding without consent face court challenges, direct action, and reputational damage. Projects with genuine consent, especially those including equity co-investment, demonstrate ``that the project has undergone a free, prior and informed consent process, meaning potential legal, reputational and social risks are significantly minimized'' (OECD/The Conversation, 2025).

\subsubsection{National Implementation Comparison}

\begin{center}
\rowcolors{2}{rowgray}{white}
\begin{tabularx}{\textwidth}{L{2.5cm}L{3cm}L{3cm}X}
\toprule
\rowcolor{secondary}
\tableheader{Country} & \tableheader{Primary Instrument} & \tableheader{FPIC Status} & \tableheader{Key Gaps} \\
\midrule
Canada & UNDRIP Act (2021); provincial variations & Legally binding FPIC (federal) & Provincial law (Ontario) not required to implement; ``duty to consult'' weaker than consent \\
USA & Federal Indian Law; NEPA; no UNDRIP adoption & Consultation only; no binding FPIC & No statutory FPIC requirement; trust responsibility unevenly enforced \\
New Zealand & Treaty of Waitangi (1840); Te Tiriti o Waitangi & Co-governance emerging; Māori data sovereignty advancing & Historic Treaty breaches not fully remediated \\
Japan & Ainu Policy Promotion Act (2019) & Cultural recognition only; no land rights, no FPIC requirement & No consultation obligation; no collective rights framework \\
Philippines & Indigenous Peoples Rights Act (1997) & First country to enshrine FPIC in national law & Implementation contested; FPIC processes subject to manipulation \\
Chile & ILO 169 ratified; Mapuche rights contested & Consultation requirement; consent contested & Mapuche land claims subject to criminal prosecution \\
\bottomrule
\end{tabularx}
\end{center}

\subsubsection{ILO Convention 169}

The ILO Indigenous and Tribal Peoples Convention (1989) is binding for ratifying states and requires consultation with Indigenous peoples when legislation or administrative measures may affect them (Article 6) and consent before relocation from lands (Article 16). It is the primary binding international instrument for countries that have not adopted UNDRIP into domestic law.

\subsection{Module B Reference: Nation-by-Nation Governance Profiles}

\textbf{Safety note:} All profiles use nation-specific names. Sacred knowledge, restricted cultural content, and confidential governance documents are excluded. Sources: government agencies, peer-reviewed research, official tribal and First Nations communications.

\subsubsection{Pacific Northwest --- Salish Sea Nations}

The central Puget Sound region and Snohomish County are part of the traditional territory of the Coast Salish peoples. The Puget Sound Regional Council recognizes nine tribal nations in the central Puget Sound region, each a sovereign nation with its own government, governing charter, and laws (PSRC, 2025). The Snohomish County Comprehensive Plan Tribal Coordination Element (2023) specifically addresses the Tulalip Tribes, the Stillaguamish Tribe of Indians, and the Sauk-Suiattle Tribe --- the three nations with trust land in Snohomish County (the Tibsfox Inc. / Paine Field prototype hub area).

\begin{center}
\rowcolors{2}{rowgray}{white}
\begin{tabularx}{\textwidth}{L{3cm}L{2.5cm}X}
\toprule
\rowcolor{secondary}
\tableheader{Nation} & \tableheader{Status / Enrollment} & \tableheader{Governance \& Economic Profile} \\
\midrule
Lummi Nation (Lhaq'temish) & Federally recognized; 5,470+ enrolled members; 13,000-acre reservation & Lummi Indian Business Council; Treaty of Point Elliott (1855); successful opposition to Gateway Pacific Terminal at Cherry Point; reef net fishing revival; Northwest Indian College \\
Tulalip Tribes & Federally recognized; trust land in Snohomish County & Successor to Snohomish Reservation; constituent tribes include Snohomish, Snoqualmie, Skykomish; major enterprises including hotel, amphitheater; significant salmon hatchery program \\
Stillaguamish Tribe & Federally recognized; 230+ members & Fish hatchery restoring Chinook and coho; Marine Stewardship and Shellfish Program; casino and hotel \\
Sauk-Suiattle Tribe & Federally recognized; trust land in Snohomish County & Land and water-based governance; Skagit watershed stewardship \\
Suquamish Tribe & Federally recognized; Port Madison Reservation & Port Madison Enterprises: second-largest private employer in Kitsap County (750+ employees) \\
Makah Indian Tribe & Federally recognized; Neah Bay & Northwesternmost point of continental U.S.; ocean and forest peoples; skilled mariners \\
\bottomrule
\end{tabularx}
\end{center}

\subsubsection{Ontario Ring of Fire --- First Nations}

\begin{center}
\rowcolors{2}{rowgray}{white}
\begin{tabularx}{\textwidth}{L{3.5cm}X}
\toprule
\rowcolor{secondary}
\tableheader{Nation / Organization} & \tableheader{Key Position} \\
\midrule
Nishnawbe Aski Nation (NAN) & Political organization representing 51 First Nation communities across Treaty 9 and Treaty 5 areas; Grand Chief Alvin Fiddler; demands FPIC as non-negotiable \\
Neskantaga First Nation & Chief Chris Moonias; Land Defence Alliance; insists on inclusion of FPIC language in all agreements; no formal position opposing development but requires consent \\
Marten Falls First Nation & Pursuing equity co-investment model; holds equity in proposed Northern Road Link infrastructure \\
15 First Nations (Regional Assessment) & Co-leading federal Regional Assessment with Impact Assessment Agency of Canada (IAAC); Terms of Reference released January 2025; 30-month timeline \\
\bottomrule
\end{tabularx}
\end{center}

\subsubsection{Japan --- Ainu of Hokkaido and Kuril Islands}

The Ainu have been officially recognized as an Indigenous people of Japan since June 2008. The Ainu Policy Promotion Act of 2019 (APPA) provides cultural recognition and tourism development subsidies but does not grant land rights, collective resource rights, or FPIC protections (Georgetown Journal of International Affairs, 2024; IWGIA, 2025). The estimated Hokkaido Ainu population is 13,118 per the 2023 government survey, though experts estimate the actual number is considerably higher.

Key governance gaps: No consultation obligation in Japanese law. Collective land rights were stripped under the Meiji Hokkaido Former Aborigines Protection Act (1899). Individual land grants were made but without collective title. The Kayano v. Hokkaido Expropriation Committee decision (1997) acknowledged the need to consider minority cultural impact in government decisions, but created no binding obligation. The Ryūkyūan (Okinawan) people are not recognized as Indigenous by the Japanese government despite UN recommendations.

\subsubsection{Aotearoa New Zealand --- Māori Governance}

The Treaty of Waitangi (1840) / Te Tiriti o Waitangi is the founding document establishing the Māori-Crown relationship. The principle of tino-rangatiratanga (Māori self-determination) is the basis for Māori data sovereignty --- governed by the CARE Principles and specifically applied through kaupapa Māori data governance frameworks. New Zealand is a global leader in Māori data sovereignty implementation, with the Crown recognizing Māori data sovereignty in government data strategy. The Whanganui River (Te Awa Tupua) was granted legal personhood in 2017, with Māori appointed as its legal guardians --- the most advanced co-governance water rights model in the Pacific Rim.

\subsection{Module C Reference: Resource Sovereignty and Critical Minerals}

\subsubsection{Canadian Ring of Fire --- James Bay Lowlands}

The Ring of Fire mineral deposit in northern Ontario is approximately 540 km northeast of Thunder Bay, within the second-largest peatland ecosystem in the world. The deposit contains chromite, cobalt, nickel, copper, and platinum --- critical minerals for electric vehicle and clean energy technology manufacturing (Canadian Mining Journal, 2025; CBC News, 2025).

Timeline of the sovereignty conflict:

\begin{center}
\rowcolors{2}{rowgray}{white}
\begin{tabularx}{\textwidth}{C{2cm}X}
\toprule
\rowcolor{primary}
\tableheader{Year} & \tableheader{Event} \\
\midrule
2014 & Ontario Liberal government releases regional framework agreement with nine Matawa chiefs; talks stall \\
2019 & Ford Progressive Conservative government abandons framework; pledges \$1 billion for infrastructure \\
2024 & Ontario still awaiting federal infrastructure funding match \\
January 2025 & IAAC and 15 First Nations release Terms of Reference for Regional Assessment \\
March 2025 & Conservative leader Poilievre pledges to ``greenlight'' all federal permits within six months; NAN Grand Chief responds: ``Our people will protect these lands, not just in the courts, but on the lands'' \\
2025--2027 & Regional Assessment 30-month timeline \\
\bottomrule
\end{tabularx}
\end{center}

The Regional Assessment is federal and non-binding on the Province of Ontario, which retains independent permitting authority. This structural gap --- federal vs. provincial jurisdiction --- is the primary mechanism by which development can proceed without meaningful consent.

\subsubsection{Pacific Geothermal --- Ring Potential}

The Ring of Fire holds over 40\% of the world's geothermal potential. By 2015, approximately 6 GW of geothermal power was installed across Ring countries (Philippines, Indonesia, New Zealand, USA). Indigenous land tenure issues are central to geothermal development: in New Zealand, Māori iwi (tribal groupings) hold surface rights and cultural authority over the Waikato volcanic plateau, where much geothermal capacity is concentrated.

\subsubsection{Fisheries and Food Sovereignty}

The 1974 Boldt Decision reaffirmed tribal fishing rights in Washington State, allocating 50\% of harvestable salmon to treaty tribes. The Lummi Nation's reef net fishing program --- the most ecologically precise salmon capture method on the coast --- is a model of food sovereignty exercised as governance. The Stillaguamish Tribe's fish hatchery restores Chinook and coho runs, an act of territorial stewardship that maintains the ecological conditions infrastructure must respect.

\subsection{Module D Reference: Cooperative Infrastructure Models}

\subsubsection{Keewaytinook Okimakanak --- KNET}

Keewaytinook Okimakanak (KO) is a tribal council of six Anishinaabe First Nations in remote northwestern Ontario. Their KNET (Kuhkenah Network) ISP is one of Canada's earliest First Nations-led internet service providers, now expanding into drone mapping, remote sensing, and AI for land monitoring. KNET is the model for Indigenous-owned digital backbone infrastructure in the Ring of Fire region.

\subsubsection{Wasaya Airways}

Wasaya Airways is owned by a coalition of First Nations and serves as a key logistics provider to Ring of Fire exploration crews. It is the clearest current example of Indigenous ownership of Ring of Fire supply chain infrastructure --- demonstrating that the co-investor model is not theoretical.

\subsubsection{Marten Falls First Nation --- Equity Co-Investment}

Marten Falls First Nation holds equity in the proposed Northern Road Link, a critical infrastructure project accessing the Ring of Fire. This equity position was secured through FPIC processes, demonstrating the principle that genuine consent produces co-ownership, not just permission.

\subsubsection{Lummi Nation --- Cherry Point Victory}

The Lummi Nation's successful opposition to the proposed Gateway Pacific Terminal coal export facility at Cherry Point, Washington, is the clearest U.S. example of treaty fishing rights as an infrastructure veto. The Lummi Nation demonstrated that treaty-protected usual and accustomed fishing grounds override port development proposals. The Lummi Indian Business Council is the governing body of the Nation under the Constitution and By-Laws approved by the federal government.

\subsubsection{Karuk Tribe --- Cultural Fire Sovereignty}

The Karuk Tribe of California/Oregon demonstrates sovereignty exercised through ecological governance. Cultural burning --- a practice governed by tribal and natural law, distinct from prescribed fire --- is asserted as a sovereign right, not a permitted activity. The Good Fire II report (2024) documents the policy recommendations for recognizing Tribal cultural fire as sovereign practice, including Tribal IRMPs (Integrated Resource Management Plans) and certification programs. The distinction between co-stewardship (collaboration without shared decision-making) and co-management (shared decision-making) is the critical governance design choice.

\subsection{Module E Reference: Indigenous Data and Digital Sovereignty}

\subsubsection{CARE Principles}

The CARE Principles for Indigenous Data Governance (Indigenous data sovereignty research community) provide the foundational framework:

\begin{center}
\rowcolors{2}{rowgray}{white}
\begin{tabularx}{\textwidth}{L{3cm}X}
\toprule
\rowcolor{secondary}
\tableheader{Principle} & \tableheader{Description} \\
\midrule
Collective Benefit & Data ecosystems shall be designed to benefit Indigenous peoples collectively, not just as data subjects \\
Authority to Control & Indigenous peoples' rights and interests in Indigenous data must be recognized, and data governance that enables this must be supported \\
Responsibility & Those working with Indigenous data have a responsibility to share how data is used \\
Ethics & Indigenous peoples' rights and wellbeing should be the primary concern at all stages of the data life cycle \\
\bottomrule
\end{tabularx}
\end{center}

\subsubsection{OCAP\textsuperscript{\textregistered} --- First Nations Information Governance}

OCAP\textsuperscript{\textregistered} (Ownership, Control, Access, Possession) is the registered standard of the First Nations Information Governance Centre (FNIGC) in Canada. It is the primary framework for First Nations data sovereignty in Canada and has direct application to Ring of Fire data governance for the Regional Assessment and infrastructure planning.

\subsubsection{Māori Data Sovereignty}

Māori data sovereignty is grounded in tino-rangatiratanga (self-determination) and Te Tiriti o Waitangi. Driven by the principle that Māori have the right to govern data about their communities, it is advancing rapidly in New Zealand's public and private sectors. Research (Oliver et al., Taylor \& Francis, 2024) identifies trust-building (whanaungatanga), self-determination (rangatiratanga), collective benefit (kotahitanga), and capacity building as core implementation requirements.

\subsubsection{Center for Tribal Digital Sovereignty}

The Center for Tribal Digital Sovereignty (NCAI / Arizona State University, launched 2024) is the first U.S. center dedicated to supporting Tribal governments in developing digital sovereignty plans. It directly addresses the gap between Indigenous data sovereignty principles and tribal government capability to implement them.

\subsection{Safety and Sensitivity Considerations}

\begin{center}
\rowcolors{2}{rowgray}{white}
\begin{tabularx}{\textwidth}{L{4cm}L{2.5cm}X}
\toprule
\rowcolor{primary}
\tableheader{Rule} & \tableheader{Boundary Type} & \tableheader{Application} \\
\midrule
Nation-specific names only & ABSOLUTE & Never ``Indigenous peoples of Canada'' as subject; always Neskantaga, Lummi, Ainu, etc. \\
Sacred knowledge excluded & ABSOLUTE & No ceremonial, restricted, or sacred content from any nation \\
OCAP\textsuperscript{\textregistered}/CARE compliance & ABSOLUTE & All data about nations must satisfy CARE Principles in how it is gathered and presented \\
Ainu naming precision & GATE & Ryūkyūans are not recognized by Japan; note carefully and without endorsing state position \\
Māori orthography & GATE & Use macrons correctly: Māori, Aotearoa, tino-rangatiratanga, iwi \\
FPIC language precision & GATE & Cite specific UNDRIP Articles; do not paraphrase as ``getting permission'' \\
Equity co-investment posture & ANNOTATE & Present as emerging model; do not imply any specific nation has agreed to any specific FIG proposal \\
Climate projections & GATE & From Pacific Islands Forum, NOAA, or peer-reviewed sources only \\
\bottomrule
\end{tabularx}
\end{center}

\subsection{Source Bibliography}

\subsubsection{Government \& Agency Sources}

\begin{itemize}
  \item United Nations, Declaration on the Rights of Indigenous Peoples (UNDRIP), 2007
  \item United Nations Permanent Forum on Indigenous Issues, Session Report, April 2025
  \item Impact Assessment Agency of Canada (IAAC), Ring of Fire Regional Assessment Terms of Reference, January 2025
  \item Snohomish County Comprehensive Plan, Tribal Coordination Element (Draft), 2023
  \item Puget Sound Regional Council (PSRC), Tribes, 2025. \url{https://www.psrc.org/about-us/tribes}
  \item Japan, Ainu Policy Promotion Act (APPA), 2019
  \item Canada, United Nations Declaration on the Rights of Indigenous Peoples Act (S.C. 2021, c. 14)
  \item Ontario, Bill 76 (UNDRIP alignment commitment), 2024
  \item Pacific Islands Forum, Pacific Regional Framework on Climate Mobility, 2024
  \item Lummi Nation, Comprehensive Economic Development Strategy (CEDS), 2023
  \item Karuk Tribe, Good Fire II Report (Executive Summary), 2024. \url{https://karuktribeclimatechangeprojects.wordpress.com}
  \item First Nations Information Governance Centre (FNIGC), OCAP\textsuperscript{\textregistered} Principles, Canada
\end{itemize}

\subsubsection{Peer-Reviewed Research}

\begin{itemize}
  \item Oliver, G., Lilley, S., Cranefield, J., \& Lewellen, M. (2024). Māori data sovereignty: contributions to data cultures in the government sector in New Zealand. \textit{The Information Society}, 40(5), 2801--2816.
  \item Karuk Tribe / Fire Research Community. (2024). Indigenous Fire Data Sovereignty: Applying Indigenous Data Sovereignty Principles to Fire Research. \textit{Fire}, 7(7), 222. MDPI.
  \item Sightline Institute. (2024). What's Misunderstood about Indigenous Cultural Fire Is Sovereignty. \url{https://www.sightline.org/2024/04/11}
  \item Georgetown Journal of International Affairs. (2024). Culture-Centered Indigenous Policies in Japan.
  \item International IDEA. (2025). Self-determination for Indigenous Peoples. \url{https://www.idea.int}
\end{itemize}

\subsubsection{Professional Organizations and News Sources}

\begin{itemize}
  \item Canadian Mining Journal. (2025, January 21). First Nations and Impact Assessment Agency finalize terms for regional assessment in Ring of Fire.
  \item Ricochet Media. (2025, April 15). `Milestone' deal reached with 15 First Nations on Ring of Fire mining development.
  \item CBC News. (2025, March 20). Poilievre criticized for pledging to fast-track Ring of Fire without Indigenous consultation.
  \item The Conversation / Phys.org. (2025, November). The path to responsible mining in northern Ontario starts with Indigenous consent.
  \item APTN News. (2025). Ontario must consult if it wants to develop Ring of Fire: Chief.
  \item Nature's Defence Canada. (2026, March). The Regional Assessment for the proposed Ring of Fire has begun.
  \item IWGIA --- International Work Group for Indigenous Affairs. (2025). The Indigenous World 2025: Japan.
  \item Ontario Ring of Fire (OROF). (2025). UNDRIP Benefits: 4 Ways Ontario First Nations Mining Tech Economy.
  \item Washington Tribes. (2025). The Tribes of Washington. \url{https://www.washingtontribes.org}
  \item Native Nations Institute, University of Arizona. Indigenous Data Sovereignty and Governance. \url{https://nni.arizona.edu}
  \item American Indian Policy Institute (AIPI) / NCAI. Center for Tribal Digital Sovereignty, launched 2024.
  \item Womble Bond Dickinson. Free, Prior, and Informed Consent \& Indigenous Consultation in Mining Development.
  \item IUCN. Resolution 8.103: Safeguarding biodiversity and human rights in energy transition mineral governance, 2025.
\end{itemize}

\newpage

%% ============================================================
%% STAGE 3 — MISSION PACKAGE
%% ============================================================
\stageheader{STAGE 3 --- MISSION PACKAGE}{tertiary}

\section{Milestone Specification: Ring of Fire --- Tribal Sovereignty \& Governance}

\subsection{Mission Objective}

Produce a comprehensive, verified research compendium on tribal sovereignty and governance across the Pacific Ring of Fire, synthesized into a GSD-ready knowledge base and FIG integration document. The compendium will provide the foundational sovereignty intelligence required for Fox Infrastructure Group's Pacific Spine partnership model --- covering legal frameworks, governance profiles, resource rights, cooperative infrastructure models, and Indigenous data sovereignty, across both the North American and Asia-Pacific arcs of the Ring.

\subsection{Architecture Overview}

\begin{asciibox}
WAVE ARCHITECTURE — Ring of Fire Sovereignty Mission

Wave 0: Foundation
    HAIKU — Shared schemas, nation profile templates,
             source index, bibliography scaffold
    [Sequential — ~2K tokens]

Wave 1: Parallel Research Tracks
    TRACK A (Sonnet): Module A (Legal) + Module C (Resources)
    TRACK B (Opus):   Module B (Governance Profiles)
    TRACK C (Sonnet): Module E (Data Sovereignty)
    [Parallel — 3 tracks — ~60K tokens total]

    CAPCOM GATE: Wave 1 review — go/no-go for Wave 2

Wave 2: Synthesis
    OPUS — Module D (Cooperative Infrastructure Models)
            Cross-module integration; FIG Sovereignty Stack mapping
    [Sequential — ~25K tokens]

    CAPCOM GATE: Synthesis review — Tibsfox approval

Wave 3: Publication + Verification
    SONNET — Document assembly, LaTeX compilation
    VERIFY  — Source validation, OCAP/CARE compliance audit
    LOG     — Commit messages, milestone summary
    [Sequential — ~15K tokens]
\end{asciibox}

\subsection{System Layers}

\begin{enumerate}
  \item \textbf{Foundation Layer} (Wave 0): Shared schemas (nation profile template, legal framework template, infrastructure model template); bibliography index; OCAP/CARE compliance checklist
  \item \textbf{Research Layer} (Wave 1): Five parallel survey documents, one per module, each self-contained and meeting minimum evidence thresholds before Wave 2 gate
  \item \textbf{Synthesis Layer} (Wave 2): Module D compendium integrating findings from Modules A, B, C, E; Sovereignty Stack $\to$ FIG mapping document
  \item \textbf{Publication Layer} (Wave 3): Final compiled XeLaTeX PDF with complete bibliography; OCAP/CARE audit log; handoff-ready GSD package
\end{enumerate}

\subsection{Deliverables}

\begin{center}
\rowcolors{2}{rowgray}{white}
\begin{tabularx}{\textwidth}{C{0.8cm}L{3.5cm}L{4cm}X}
\toprule
\rowcolor{tertiary}
\tableheader{\#} & \tableheader{Deliverable} & \tableheader{Acceptance Criteria} & \tableheader{Spec} \\
\midrule
D1 & Legal Frameworks Survey & All 6 countries covered; all FPIC citations use UNDRIP Article numbers & Module A document \\
D2 & Nation Governance Profiles & 12+ profiles; nation-specific names; OCAP/CARE compliant & Module B document \\
D3 & Resource Sovereignty Compendium & Ring of Fire (Ontario) case complete; geothermal + fisheries sections & Module C document \\
D4 & Data Sovereignty Survey & CARE, OCAP, Māori, USIDSN, CTDS covered & Module E document \\
D5 & Cooperative Infrastructure Synthesis & 5 case studies; FIG Sovereignty Stack mapping & Module D document \\
D6 & Compiled Research PDF & XeLaTeX 3-pass; zero errors; 25+ sources & Final PDF \\
D7 & OCAP/CARE Audit Log & Zero generic identifiers; zero sacred content violations & Audit document \\
\bottomrule
\end{tabularx}
\end{center}

\subsection{Component Breakdown}

\begin{center}
\rowcolors{2}{rowgray}{white}
\begin{tabularx}{\textwidth}{L{3.5cm}L{3cm}C{1.5cm}C{2cm}}
\toprule
\rowcolor{primary}
\tableheader{Component} & \tableheader{Dependencies} & \tableheader{Model} & \tableheader{Est. Tokens} \\
\midrule
Wave 0: Schemas \& Templates & None & Haiku & 2,000 \\
Module A: Legal Frameworks & Wave 0 & Sonnet & 15,000 \\
Module B: Governance Profiles & Wave 0 & Opus & 20,000 \\
Module C: Resource Sovereignty & Wave 0 & Sonnet & 12,000 \\
Module E: Data Sovereignty & Wave 0 & Sonnet & 10,000 \\
Module D: Cooperative Synthesis & A, B, C, E & Opus & 25,000 \\
Wave 3: Assembly \& Verification & A--E, D & Sonnet & 15,000 \\
\textbf{Total (estimated)} & & & \textbf{99,000} \\
\bottomrule
\end{tabularx}
\end{center}

\subsection{Model Assignment Rationale}

\begin{center}
\rowcolors{2}{rowgray}{white}
\begin{tabularx}{\textwidth}{L{2cm}L{2cm}X}
\toprule
\rowcolor{primary}
\tableheader{Model} & \tableheader{Allocation} & \tableheader{Rationale} \\
\midrule
Opus & ~30\% & Cultural sensitivity judgment (Module B); cross-module synthesis requiring pattern extraction across five bodies of research (Module D); OCAP/CARE compliance decisions \\
Sonnet & ~60\% & Structured survey compilation (Modules A, C, E); document assembly; source validation; XeLaTeX publication \\
Haiku & ~10\% & Scaffolding, shared schemas, bibliography templates, nation profile stubs \\
\bottomrule
\end{tabularx}
\end{center}

\section{Wave Execution Plan}

\subsection{Wave Summary}

\begin{center}
\rowcolors{2}{rowgray}{white}
\begin{tabularx}{\textwidth}{C{1.2cm}L{3cm}C{2cm}C{2cm}C{2cm}}
\toprule
\rowcolor{primary}
\tableheader{Wave} & \tableheader{Work} & \tableheader{Mode} & \tableheader{Models} & \tableheader{Gate} \\
\midrule
0 & Foundation (schemas, templates, source index) & Sequential & Haiku & Auto \\
1 & Modules A + B + C + E (parallel) & Parallel & Sonnet + Opus & CAPCOM \\
2 & Module D synthesis + FIG mapping & Sequential & Opus & CAPCOM \\
3 & Assembly, verification, publication & Sequential & Sonnet & Auto \\
\bottomrule
\end{tabularx}
\end{center}

\subsection{Wave 0: Foundation}

\begin{center}
\rowcolors{2}{rowgray}{white}
\begin{tabularx}{\textwidth}{L{4cm}XC{2cm}}
\toprule
\rowcolor{primary}
\tableheader{Task} & \tableheader{Output} & \tableheader{Agent} \\
\midrule
Nation profile template & Standard schema: name, governance structure, treaty/legal status, economic posture, known infrastructure positions & Haiku \\
Legal framework template & Standard schema: instrument name, date, parties, FPIC provision, national implementation status & Haiku \\
Infrastructure model template & Standard schema: name, nation(s), type, outcome, replicability signal & Haiku \\
Source index & Master bibliography stub with all known sources from this document, ready for agent annotation & Haiku \\
OCAP/CARE compliance checklist & Audit checklist: nation-specific naming, no generic identifiers, no sacred content, CARE principles per section & Haiku \\
\bottomrule
\end{tabularx}
\end{center}

\subsection{Wave 1: Parallel Research Tracks}

\textbf{Track A --- Sonnet:}

\begin{center}
\rowcolors{2}{rowgray}{white}
\begin{tabularx}{\textwidth}{L{3cm}XC{2cm}}
\toprule
\rowcolor{secondary}
\tableheader{Task} & \tableheader{Output} & \tableheader{Agent} \\
\midrule
Module A: Legal survey & UNDRIP, ILO 169, six-country national implementation table & Sonnet \\
Module C: Resources & Ring of Fire (Ontario) case; geothermal map; fisheries sovereignty & Sonnet \\
\bottomrule
\end{tabularx}
\end{center}

\textbf{Track B --- Opus:}

\begin{center}
\rowcolors{2}{rowgray}{white}
\begin{tabularx}{\textwidth}{L{3cm}XC{2cm}}
\toprule
\rowcolor{secondary}
\tableheader{Task} & \tableheader{Output} & \tableheader{Agent} \\
\midrule
Module B: Governance profiles & 12+ nation profiles using nation-specific names; OCAP/CARE compliant & Opus \\
\bottomrule
\end{tabularx}
\end{center}

\textbf{Track C --- Sonnet:}

\begin{center}
\rowcolors{2}{rowgray}{white}
\begin{tabularx}{\textwidth}{L{3cm}XC{2cm}}
\toprule
\rowcolor{secondary}
\tableheader{Task} & \tableheader{Output} & \tableheader{Agent} \\
\midrule
Module E: Data sovereignty & CARE, OCAP, Māori, USIDSN, CTDS; Indigenous digital infrastructure cases & Sonnet \\
\bottomrule
\end{tabularx}
\end{center}

\subsection{Wave 2: Synthesis}

\begin{center}
\rowcolors{2}{rowgray}{white}
\begin{tabularx}{\textwidth}{L{4cm}XC{2cm}}
\toprule
\rowcolor{primary}
\tableheader{Task} & \tableheader{Output} & \tableheader{Agent} \\
\midrule
Module D: Cooperative infrastructure synthesis & Five case studies with cross-module integration; replicability signals & Opus \\
FIG Sovereignty Stack mapping & Five-layer Sovereignty Stack mapped to FIG co-op subsidiaries (FoxFiber, FoxCompute, FoxData, FoxEnergy, FoxEdu); partnership entry points per layer & Opus \\
Cross-module dependency validation & Confirm all D inputs are present and sufficient from Modules A, B, C, E & INTEG \\
\bottomrule
\end{tabularx}
\end{center}

\subsection{Wave 3: Publication and Verification}

\begin{center}
\rowcolors{2}{rowgray}{white}
\begin{tabularx}{\textwidth}{L{4cm}XC{2cm}}
\toprule
\rowcolor{primary}
\tableheader{Task} & \tableheader{Output} & \tableheader{Agent} \\
\midrule
Document assembly & LaTeX compilation of all modules into final PDF (3-pass XeLaTeX) & Sonnet \\
Source validation & Verify all 25+ sources are government/peer-reviewed; no blogs, no entertainment media & VERIFY \\
OCAP/CARE audit & Zero generic identifiers; zero sacred content; CARE principles per section & VERIFY \\
Commit messages + milestone summary & Audit trail; v1.50 milestone integration log & LOG \\
\bottomrule
\end{tabularx}
\end{center}

\subsection{Cache Optimization Strategy}

\begin{itemize}
  \item Modules A and C share legal framework context --- compile legal source index first in Wave 0 and pass as shared cache to both Track A agents
  \item Module B nation profiles use nation profile template from Wave 0 as a structured prompt prefix --- pre-load template into all Track B agent contexts
  \item Module D synthesis agent receives compressed summaries (not full text) of Modules A, B, C, E --- use 1,000-token executive summary per module to preserve context window for synthesis reasoning
  \item Wave 3 assembly agent works from the Module D synthesis document as primary input; Modules A--E are secondary references fetched on demand
  \item FIG mapping section in Module D is time-sensitive for Tibsfox review --- flag as CAPCOM priority at Wave 2 gate
\end{itemize}

\subsection{Token Budget}

\begin{center}
\rowcolors{2}{rowgray}{white}
\begin{tabularx}{\textwidth}{L{4cm}C{3cm}C{3cm}}
\toprule
\rowcolor{primary}
\tableheader{Component} & \tableheader{Estimated Tokens} & \tableheader{Model} \\
\midrule
Wave 0: Foundation & 2,000 & Haiku \\
Module A: Legal frameworks & 15,000 & Sonnet \\
Module B: Governance profiles & 20,000 & Opus \\
Module C: Resource sovereignty & 12,000 & Sonnet \\
Module E: Data sovereignty & 10,000 & Sonnet \\
Module D: Synthesis + FIG mapping & 25,000 & Opus \\
Wave 3: Assembly + verification & 15,000 & Sonnet \\
\midrule
\textbf{Total estimated} & \textbf{99,000} & Mixed \\
\bottomrule
\end{tabularx}
\end{center}

\subsection{Dependency Graph}

\begin{asciibox}
DEPENDENCY GRAPH

Wave 0: [W0-Schema] ─────────────────────────────────┐
                                                      │
Wave 1:  [A: Legal] ──────────────────────┐           │
         [B: Governance] ─────────────────┼───►[D: Synthesis]
         [C: Resources] ──────────────────┤           │
         [E: Data Sov.] ──────────────────┘           │
                              │                        │
                              ▼                        │
           [INTEG: Cross-Module Validation]            │
                              │                        │
                              ▼                        ▼
Wave 3:          [VERIFY] ◄──────────── [Assembly ◄─ W0-Schema]
                    │
                    ▼
                [LOG ─── Final PDF + Audit]

Critical path: B (Governance Profiles) — longest track, Opus-heavy
               D (Synthesis) — depends on all four Wave 1 outputs
\end{asciibox}

\section{Test Plan}

\subsection{Test Categories}

\begin{center}
\rowcolors{2}{rowgray}{white}
\begin{tabularx}{\textwidth}{L{3.5cm}C{1.5cm}L{2cm}X}
\toprule
\rowcolor{primary}
\tableheader{Category} & \tableheader{Count} & \tableheader{Priority} & \tableheader{Failure Action} \\
\midrule
Safety-Critical (OCAP/CARE) & 6 & Mandatory-pass & BLOCK mission until resolved \\
Core Functionality & 8 & High & Flag and repair before Wave 3 \\
Integration & 4 & Medium & Note and document; proceed with caveat \\
Edge Cases & 3 & Low & Document for v2.0 backlog \\
\bottomrule
\end{tabularx}
\end{center}

\subsection{Safety-Critical Tests}

\begin{center}
\rowcolors{2}{rowgray}{white}
\begin{tabularx}{\textwidth}{C{1.2cm}L{3.5cm}X}
\toprule
\rowcolor{primary}
\tableheader{Test ID} & \tableheader{Verifies} & \tableheader{Expected Result} \\
\midrule
SC-01 & No generic Indigenous identifiers & Zero instances of ``Indigenous communities'' or similar without nation-specific attribution in Module B \\
SC-02 & No sacred or restricted content & Zero reproduction of ceremonial, sacred, or restricted content from any nation \\
SC-03 & FPIC language uses UNDRIP Articles & All FPIC references include specific Article citations (10, 19, 28, or 32); zero informal paraphrase \\
SC-04 & Ainu/Ryūkyūan distinction preserved & Ryūkyūans noted as not recognized by Japan without endorsing state position \\
SC-05 & Māori orthography correct & All Māori terms use correct macrons; no anglicized alternatives \\
SC-06 & Source quality gate & Zero sources from entertainment media, blogs, or unsourced claims; 100\% government/peer-reviewed/professional \\
\bottomrule
\end{tabularx}
\end{center}

\subsection{Verification Matrix}

\begin{center}
\rowcolors{2}{rowgray}{white}
\begin{tabularx}{\textwidth}{L{0.8cm}X L{3.5cm}}
\toprule
\rowcolor{primary}
\tableheader{SC} & \tableheader{Success Criterion} & \tableheader{Test IDs} \\
\midrule
SC1 & Legal framework survey covering 6+ countries & SC-03, CORE-01, CORE-02 \\
SC2 & 12+ nation profiles, nation-specific, OCAP compliant & SC-01, SC-02, SC-04, SC-05 \\
SC3 & Canadian Ring of Fire case study complete & CORE-03, CORE-04 \\
SC4 & 5 cooperative infrastructure cases documented & CORE-05, INTEG-01 \\
SC5 & Data sovereignty landscape survey & CORE-06, INTEG-02 \\
SC6 & FIG Sovereignty Stack mapping & INTEG-01, INTEG-02 \\
SC7 & OCAP/CARE compliance audit passed & SC-01, SC-02, SC-06 \\
SC8 & Module dependency graph verified & INTEG-03, INTEG-04 \\
SC9 & 25+ verified sources & SC-06, CORE-07 \\
SC10 & XeLaTeX compiles clean, 3-pass & CORE-08 \\
SC11 & FPIC language uses UNDRIP Articles & SC-03 \\
SC12 & Zero additional context needed for handoff & INTEG-04 \\
\bottomrule
\end{tabularx}
\end{center}

\section{Mission Crew Manifest}

\begin{center}
\rowcolors{2}{rowgray}{white}
\begin{tabularx}{\textwidth}{L{2cm}L{2.5cm}X}
\toprule
\rowcolor{primary}
\tableheader{Role} & \tableheader{Agent} & \tableheader{Responsibility} \\
\midrule
FLIGHT & Orchestrator & Mission direction; go/no-go gates; CAPCOM escalation \\
PLAN & Planner & Wave decomposition; dependency management; parallel track optimization \\
SCOUT & Research Agent & Source gathering; web research; agency document retrieval \\
EXEC-A & Executor Alpha & Module A (Legal) + Module C (Resources) production \\
EXEC-B & Executor Beta & Module E (Data Sovereignty) production + Wave 3 assembly \\
CRAFT-GOV & Governance Specialist & Module B: nation-by-nation governance profiles (Opus) \\
CRAFT-INFRA & Infrastructure Specialist & Module D: cooperative infrastructure synthesis (Opus) \\
VERIFY & Verification Agent & Source validation; OCAP/CARE compliance audit; SC-06 gate \\
INTEG & Integration Agent & Cross-module dependency validation; FIG mapping coherence check \\
CAPCOM & HITL Interface & Human approval at Wave 1 and Wave 2 gates; only Tibsfox-facing channel \\
BUDGET & Resource Manager & Token budget tracking; session planning; wave boundary alerts \\
LOG & Chronicle Agent & Commit messages; audit trail; milestone summary for v1.50 \\
\bottomrule
\end{tabularx}
\end{center}

\section{Execution Summary}

\begin{center}
\rowcolors{2}{rowgray}{white}
\begin{tabularx}{\textwidth}{L{4cm}X}
\toprule
\rowcolor{primary}
\tableheader{Metric} & \tableheader{Value} \\
\midrule
Activation Profile & Squadron (12 roles) \\
Total Modules & 5 (A: Legal, B: Governance, C: Resources, D: Synthesis, E: Data) \\
Parallel Tracks (Wave 1) & 3 (A+C, B, E) \\
Critical Path & Module B (Governance, Opus) $\to$ Module D (Synthesis, Opus) \\
CAPCOM Gates & 2 (after Wave 1, after Wave 2) \\
Estimated Token Budget & 99,000 \\
Context Windows & 18--24 across 4--6 sessions \\
Primary Sensitivity Framework & OCAP\textsuperscript{\textregistered} / CARE / UNDRIP \\
Safety-Critical Tests & 6 (all BLOCK-authority) \\
Target Deliverable & Compiled XeLaTeX PDF + FIG Sovereignty Stack mapping \\
Integration Target & Fox Infrastructure Group Pacific Spine partnership architecture \\
Mission Package Status & Ready for GSD Orchestrator handoff \\
\bottomrule
\end{tabularx}
\end{center}

\vspace{2em}

\begin{quotebox}
\textit{``We've lived like that for a long time. It needs to be done in a way where we are at the table, not on the menu. And we need to be part of the discussion.''}

\vspace{0.5em}
--- Sol Mamakwa, Kiiwetinoong MPP (Anishinaabe / Oji-Cree), responding to proposals to fast-track the Ring of Fire without Indigenous consultation, March 2025

\vspace{1em}
The Ring of Fire is not an engineering problem. It is a governance problem that engineering can only address after the governance question is settled. The Sovereignty Stack is the correct order of operations. Build from Layer 1. Everything else follows.

\textit{--- Tibsfox Inc. / Fox Infrastructure Group Design Principle}
\end{quotebox}

\end{document}
